\section{Cut elimination for the exact intuitionistic sheet}
\label{sec:p07-annex-intuitionistic-cut-elimination}

\begin{annexitem}{annex:p7.lj.alpha-substitution}{thm:p07-lj-alpha-substitution}{Substitution and fresh eigenparameters}
\begin{proof}
First record the binding facts needed for the proof representation.  Term
substitution respects alpha-equivalence.  A simultaneous finite substitution
therefore acts on alpha-classes.  It leaves the tree of logical connectives
and quantifiers unchanged, hence preserves \(\rho\).  In particular, if
\(z\) is fresh for \(B\) and \(t\), then
\[
             [(B[z/x])[t/z]]_\alpha=[B[t/x]]_\alpha.
\]
More generally, substitution commutes with a quantifier-rule instance
after its binder and eigenparameter have been made fresh for the finite
substitution's domain and the free variables in its range.

To justify the use of alpha-classes, first observe that each generating
bound renaming preserves the free-variable set and the logical constructor
tree.  It also preserves which quantifier governs each bound occurrence:
the replacement variable is absent from the matrix before renaming.
Reflexivity, symmetry, transitivity and constructor congruence preserve
these three properties.  Given finitely many formulas being compared,
assign a distinct fresh variable to each quantifier position of their
common constructor tree, avoiding all their free variables and all terms
involved, and choose them outside every old bound-variable name as well.
Rename the binders from the outside in to those assigned variables.
The definition of alpha-equivalence licenses each step.  Alpha-equivalent
formulas now have identical free tokens and identical bound tokens at
corresponding positions, hence identical resulting literal representatives.
Applying a capture-avoiding substitution to this common representative,
using the same fresh bound names, gives the same result on either side.
Changing those fresh names is itself a sequence of generating bound
renamings.  This proves substitution congruence and independence of the
chosen fresh representatives.  It also shows that quantified alpha-equivalent
formulas, opened at a common fresh variable, have alpha-equivalent matrices.

Here is a structural proof of these facts that does not assume identities
between least-fresh literal strings.  In an atomic formula, substitute
recursively in its terms; variable, constant and function-application
cases give the term substitution equations.  The bottom case is unchanged.
For a binary connective apply the induction hypothesis to both children.
At \(QxB\), choose a binder \(w\) outside the finitely many variables
of the formulas and terms involved.  Represent the class as
\(Qw(B[w/x])\).  All substitutions being compared now act on its matrix
without capturing \(w\).  The matrix induction hypothesis proves the
desired equality of classes, and putting \(Qw\) back proves the quantified
case.  Different fresh choices are alpha-equivalent by the same argument.
This also proves that alpha-equivalent quantified formulas have
alpha-equivalent instances at any fixed term.  No term substitution adds
or removes a logical symbol, which proves the rank assertion.

For the weak quantifier cases below, the precise composition identity is
\[
 [(B[t/w])\theta]_\alpha
       =[(B\theta)[t\theta/w]]_\alpha,
\]
provided \(\theta\) fixes \(w\) and no term in its finite nontrivial
range contains \(w\).  Term induction proves this identity for variables
and function applications: the variable \(w\) gives \(t\theta\),
and any other variable gives its \(\theta\)-image, which contains no
\(w\).  Atomic formulas apply those term identities.  Bottom is unchanged
and the connective cases commute componentwise.  At a quantifier, choose
its binder outside the finitely many variables being compared and apply
the matrix induction hypothesis.  The resulting equality is on alpha-classes,
as required.  Before using this identity on an instance of \(QxB\),
choose its representative's binder \(w\) fresh for the substitution's
domain and range; this satisfies the two stated conditions.

Projection of a literal proof replaces every formula occurrence by its
class and changes neither the tree nor the rule instances.  It is thus
height preserving.  We prove height-preserving substitution for cut-free
alpha-proofs by induction on tree height, uniformly for every finite
simultaneous substitution \(\theta\).  An identity remains an identity
of classes; a bottom-left initial remains bottom left.  WL and CL retain
their selected context occurrences.  Each propositional rule retains its
shape: contexts shared by \(\land R\) and \(\lor L\) receive the same
substitution on both branches, and the two contexts of \(\to L\) retain
their split.  Apply the induction hypothesis to the premises and reapply
the rule.

For \(\forall L\) and \(\exists R\), choose a fresh representative
of the quantified binder.  After applying \(\theta\), the active
instance is the instance of the substituted matrix at the substituted
witness term, by the binding facts just proved.  Thus the same weak
quantifier rule applies to the substituted premise proof.

For \(\forall R\) or \(\exists L\), let its old eigenparameter be
\(y\).  Choose \(z\) outside the variables of the rule instance, the
domain and term ranges of \(\theta\), and any additional finite
forbidden set.  Apply the induction hypothesis to the strictly smaller
premise proof with substitution \([z/y]\).  Its height does not increase.
Because \(y\) is absent from the free variables of the conclusion,
this renaming changes no conclusion class.  The fresh-instance identity
identifies the new active premise as the instance at \(z\).  Apply the
same height induction hypothesis to this still shorter proof with
\(\theta\).  The eigenparameter \(z\) cannot occur freely in the
substituted conclusion: it is absent both from the old conclusion and
from the substitution range.  Reapply the eigenvariable rule.  The
result has height at most that of the original inference.

This argument also supplies eigenparameter freshening itself.  At each
eigenvariable inference perform its parameter change on the premise
subtree, where the preceding induction applies; then proceed to the
remaining finitely many inferences, adding chosen parameters to the
forbidden set.  Previously fixed outside labels are unchanged by the
local eigencondition.  This obtains distinct fresh parameters where
needed and avoids any prescribed finite set, without increasing height.
There is no circular call at the height of the inference being treated.
\end{proof}
\end{annexitem}

\begin{annexitem}{annex:p7.lj.context-collapse}{thm:p07-lj-context-collapse}{Collapse of supplied contexts}
\begin{proof}
For every formula \(B\) with multiplicity \(g_B\) in \(\Gamma\),
the designated \(m\) copies contain \(mg_B\) occurrences of \(B\).
Apply CL exactly \((m-1)g_B\) times to pairs of designated occurrences,
leaving \(g_B\) of them.  Do this for every formula in the finite support
of \(\Gamma\).  The antecedent left over is exactly
\(\Gamma,\Pi\).  Occurrences from \(\Pi\) are not chosen for these
contractions, even if they carry the same formula.  All the steps are
registered CL steps and introduce neither Cut nor Mix.  For empty
\(\Gamma\), or for \(m=1\), no step is necessary.  The argument uses
only multiset multiplicities and so applies to both proof representations.
\end{proof}
\end{annexitem}

\begin{annexitem}{annex:p7.lj.bottom-replacement}{thm:p07-lj-bottom-replacement}{Bottom-succedent replacement without right weakening}
\begin{proof}
Induct on a cut-free alpha-proof of \(\Gamma\Rightarrow\bot\).
An identity leaf must be \(\bot\Rightarrow\bot\); replace it by
bottom left with succedent \(C\).  A bottom-left leaf likewise remains
a bottom-left leaf after changing its succedent.

For WL and CL change the succedent in the premise by induction and
restore the structural step.  For \(\land L\), change the succedent
in its premise \(B,F,\Gamma_0\Rightarrow\bot\), then restore the
conjunction.  For \(\lor L\), make the change in both premises;
their common context is unchanged, so the rule applies again.  For
\(\to L\), leave the first premise \(\Sigma\Rightarrow B\)
unchanged and apply induction only to the second premise
\(F,\Theta\Rightarrow\bot\); restore the implication-left rule
with succedent \(C\).

For \(\forall L\), change the succedent in its premise and restore
the same witness instance.  For \(\exists L\), first use
\ref{thm:p07-lj-alpha-substitution} to freshen the eigenparameter away
from \(C\).  The freshened premise has smaller height than the original
proof.  Apply induction to it and restore \(\exists L\).  Its
eigenparameter is absent from the free variables of the new conclusion.

No right logical introduction can have primitive \(\bot\) as its
succedent.  There are no right structural rules, and the input has no
Cut or Mix.  These are all cases in \ref{sheet:lj}.
\end{proof}
\end{annexitem}

\begin{annexitem}{annex:p7.lj.marked-multicut}{thm:p07-lj-marked-multicut}{Complete marked-multicut induction}
\begin{proof}
All formulas in this item are alpha-classes, written with formula letters
for readability.  The input proofs contain neither Cut nor Mix.  If
\(n=0\), add the finitely many occurrences of \(\Gamma\) to \(E\)
by WL and finish.  For \(n>0\), induct lexicographically on
\(\bigl(\rho(A),h(D)+h(E)\bigr)\).  Thus a smaller-rank call may
have arbitrary finite proof heights, whereas a same-rank call must have
strictly smaller summed input height.  We always keep unselected
occurrences, including occurrences equal to \(A\).

\paragraph{Initial cases.}
If \(D\) is identity \(A\Rightarrow A\), contract the \(n\)
selected occurrences of \(A\) in \(E\) to one, leaving \(\Pi\)
untouched.  The result is \(A,\Pi\Rightarrow C\), as required.
If \(D\) is bottom left, \(\Gamma\) contains a bottom occurrence,
and bottom left directly gives \(\Gamma,\Pi\Rightarrow C\).
If \(E\) is identity, its uncontexted antecedent has just one
occurrence.  Since \(n>0\), that occurrence is the selected \(A\),
\(\Pi\) is empty and \(C=A\).  Return \(D\).
If \(E\) is bottom left and an unselected bottom remains in
\(\Pi\), close the target directly by bottom left.  Otherwise a
selected occurrence is bottom, so \(A=\bot\).  Apply
\ref{thm:p07-lj-bottom-replacement} to \(D\), obtaining
\(\Gamma\Rightarrow C\), and weaken \(\Pi\) into it.
This case does not assume that \(\Gamma\) itself contains bottom.

\paragraph{Nonprincipal last inferences of the left proof.}
The selected left occurrence is its sole succedent \(A\).  Thus its
last inference is nonprincipal exactly when it is a structural or logical
left inference.  In each case below the recursive call uses an immediate
premise of \(D\) and the original \(E\), so the height sum decreases.

For WL, \(D\) has premise \(\Gamma_0\Rightarrow A\) and
conclusion \(B,\Gamma_0\Rightarrow A\).  Induction gives
\(\Gamma_0,\Pi\Rightarrow C\); weaken \(B\) back in.
For CL, its premise is \(B,B,\Gamma_0\Rightarrow A\).  Induction
gives \(B,B,\Gamma_0,\Pi\Rightarrow C\); contract the two
designated copies back to one.  In both cases \(B\) may equal \(A\):
these are occurrences in the supplied left context, not right-side marks.

For \(\land L\), apply induction to
\(B_1,B_2,\Gamma_0\Rightarrow A\), then restore the left
conjunction in \(B_1,B_2,\Gamma_0,\Pi\Rightarrow C\).
For \(\lor L\), apply induction separately to
\(B_1,\Gamma_0\Rightarrow A\) and
\(B_2,\Gamma_0\Rightarrow A\).  The resulting premises have the
same side context \(\Gamma_0,\Pi\), so \(\lor L\) gives
\(B_1\lor B_2,\Gamma_0,\Pi\Rightarrow C\).

For \(\to L\), the premises of \(D\) are
\(\Sigma\Rightarrow B\) and \(F,\Theta\Rightarrow A\).
Keep the first unchanged and apply induction to the second and \(E\).
Its output is \(F,\Theta,\Pi\Rightarrow C\).  Reapplying the
split-context rule gives
\(B\to F,\Sigma,\Theta,\Pi\Rightarrow C\), the required
antecedent \(\Gamma,\Pi\).  There is no cut on the first premise's
succedent, even if it happens to be \(A\).

For \(\forall L\), induction on its premise
\(B[t/x],\Gamma_0\Rightarrow A\) gives
\(B[t/x],\Gamma_0,\Pi\Rightarrow C\).  Restore \(\forall L\)
with the same term \(t\).  For \(\exists L\), first freshen its
eigenparameter to \(z\), outside the variables of \(E\) and hence
outside \(\Pi,C\).  Its premise is
\(B[z/x],\Gamma_0\Rightarrow A\).  Induction gives
\(B[z/x],\Gamma_0,\Pi\Rightarrow C\), and \(\exists L\)
applies because \(z\) is fresh for the new conclusion.  Freshening
does not increase the premise height.

This covers every left nonprincipal case.  If none applies and \(D\)
is not an initial sequent, its last inference is a right introduction
of \(A\).

\paragraph{Structural last inferences of the right proof.}
If \(E\) ends in WL introducing an unselected \(B\), transport
the marks to the premise, apply induction and weaken \(B\) back in.
If the introduced occurrence is selected, it is \(A\).  Transport
only the other \(n-1\) marks to the premise and apply induction, or
the zero-mark padding case when \(n=1\).  The result is already
\(\Gamma,\Pi\Rightarrow C\); the selected weakening is omitted.

If \(E\) ends in CL contracting a selected occurrence, select both
of its predecessor occurrences.  The premise has \(n+1\) marks,
and smaller height.  Induction gives the desired result directly.
If the contracted occurrence is unselected, leave both predecessor
occurrences unselected.  Induction produces
\(\Gamma,B,B,\Pi_0\Rightarrow C\); CL gives
\(\Gamma,B,\Pi_0\Rightarrow C\).  Again, an unselected \(B\)
may equal \(A\) without being deleted.

\paragraph{Nonprincipal logical last inferences of the right proof.}
For a logical left inference in this paragraph its principal occurrence
is unselected; its exposed constituent occurrences are therefore
unselected.  For a logical right inference every selected antecedent
occurrence is a context occurrence.  Each recursive call uses the
original \(D\) and an immediate premise of \(E\), after at most
height-preserving freshening.  A branch with no marks uses padding.

For \(\land L\), the premise is
\(B,F,A^n,\Pi_0\Rightarrow C\).  Induction gives
\(B,F,\Gamma,\Pi_0\Rightarrow C\); restore the conjunction.
For \(\land R\), apply induction to both premises
\(A^n,\Pi\Rightarrow B\) and \(A^n,\Pi\Rightarrow F\).
Both output contexts are exactly \(\Gamma,\Pi\), so the shared-context
rule applies with just one copy of that context in its conclusion.

For an unselected \(\lor L\), the premises are
\(B,A^n,\Pi_0\Rightarrow C\) and
\(F,A^n,\Pi_0\Rightarrow C\).  Apply induction to both and
restore \(\lor L\) with common side context \(\Gamma,\Pi_0\).
For either \(\lor R_1\) or \(\lor R_2\), apply induction to
its one premise and restore that same injection.

For an unselected \(\to L\), write its premises as
\(\Sigma\Rightarrow B\) and \(F,\Theta\Rightarrow C\).
The marks split into \(n_1\) occurrences in \(\Sigma\) and
\(n_2\) in \(\Theta\), with \(n_1+n_2=n\).  Let
\(\Sigma_0,\Theta_0\) be their unselected remainders.  Induction
or zero-mark padding gives
\[
 \Gamma,\Sigma_0\Rightarrow B,
 \qquad F,\Gamma,\Theta_0\Rightarrow C.
\]
The active \(F\) is unselected even if it equals \(A\).
The split-context implication rule now gives
\(B\to F,\Gamma,\Sigma_0,\Gamma,\Theta_0\Rightarrow C\).
Collapse the two supplied copies of \(\Gamma\) by
\ref{thm:p07-lj-context-collapse}.  This gives exactly the required
\(\Gamma,B\to F,\Sigma_0,\Theta_0\Rightarrow C\).

For \(\to R\), its premise is \(B,A^n,\Pi\Rightarrow F\).
The active assumption \(B\), to be discharged, is unselected even
when it equals \(A\).  Induction gives
\(B,\Gamma,\Pi\Rightarrow F\), and \(\to R\) discharges
that occurrence, giving \(\Gamma,\Pi\Rightarrow B\to F\).

For an unselected \(\forall L\), apply induction to
\(B[t/x],A^n,\Pi_0\Rightarrow C\) and restore the same witness
term.  For \(\exists R\), apply induction to
\(A^n,\Pi\Rightarrow B[t/x]\) and restore the same existential
introduction.  Neither rule has an eigenvariable condition on this sheet.

For \(\forall R\), freshen its eigenparameter to \(z\), outside
the variables of \(D\), as well as its own conclusion.  Induction
on \(A^n,\Pi\Rightarrow B[z/x]\) yields
\(\Gamma,\Pi\Rightarrow B[z/x]\).  Since \(z\) is not free
in \(\Gamma\), the new conclusion satisfies the eigencondition and
\(\forall R\) applies.  For an unselected \(\exists L\), make
the same freshness choice before applying induction to
\(B[z/x],A^n,\Pi_0\Rightarrow C\).  Its output is
\(B[z/x],\Gamma,\Pi_0\Rightarrow C\), and the fresh parameter
allows \(\exists L\) to be restored.

\paragraph{Principal conjunction.}
It remains that \(D\) introduces \(A\) on the right and \(E\)
introduces a selected \(A\) on the left.  In each of the following
cases first remove the residual \(n-1\) selected occurrences from
the right premises using the original \(D\).  These calls have
smaller original right-premise height; for no residual marks they are
just padding.  Subsequent calls have smaller rank.

Let \(A=B\land F\).  The left premises are
\(D_B:\Gamma\Rightarrow B\) and
\(D_F:\Gamma\Rightarrow F\).  The right premise is
\(E_0:B,F,A^{n-1},\Pi\Rightarrow C\).  Residual cleanup yields
\(B,F,\Gamma,\Pi\Rightarrow C\).  Use the smaller-rank lemma
on the displayed \(B\), with \(D_B\), to obtain
\(\Gamma,F,\Gamma,\Pi\Rightarrow C\); collapse its two
supplied contexts.  Use the smaller-rank lemma on the displayed \(F\),
with \(D_F\), to obtain \(\Gamma,\Gamma,\Pi\Rightarrow C\),
and collapse again.  Both constituent ranks are strictly less than
\(\rho(A)\), so the heights of these new proofs are immaterial.

\paragraph{Principal disjunction, both injections.}
Let \(A=B\lor F\).  If \(D\) ends with \(\lor R_1\), let
its premise be \(D_B:\Gamma\Rightarrow B\).  Choose the
\(B\)-premise \(B,A^{n-1},\Pi\Rightarrow C\) of the final
\(\lor L\) in \(E\).  Residual cleanup gives
\(B,\Gamma,\Pi\Rightarrow C\).  Apply the smaller-rank lemma
on that displayed \(B\), with \(D_B\), then collapse the two
copies of \(\Gamma\).  If \(D\) ends with \(\lor R_2\), use
its premise \(\Gamma\Rightarrow F\) and the \(F\)-premise
of \(E\) instead.  Cleanup, the smaller-rank operation on \(F\),
and context collapse give the same target.  Only the branch chosen by
the right introduction is needed in the resulting derivation.

\paragraph{Principal implication.}
Let \(A=B\to F\).  The premise of \(D\) is
\(D_0:B,\Gamma\Rightarrow F\).  The two premises of \(E\)
are \(E_1:\Sigma\Rightarrow B\) and
\(E_2:F,\Theta\Rightarrow C\).  Residual marks are distributed
between \(\Sigma\) and \(\Theta\); write their unselected
remainders as \(\Sigma_0\) and \(\Theta_0\).  Thus
\(\Pi=\Sigma_0,\Theta_0\).  Cleanup with the original \(D\)
and these two original premise proofs gives
\[
 R_1:\Gamma,\Sigma_0\Rightarrow B,
 \qquad R_2:F,\Gamma,\Theta_0\Rightarrow C.
\]
Use padding in a branch without residual marks.  Apply the smaller-rank
lemma on \(B\) to \(R_1\) and \(D_0\), selecting only the
displayed active \(B\) of \(D_0\).  The result is
\(\Gamma,\Sigma_0,\Gamma\Rightarrow F\).  Collapse the two
supplied contexts to obtain \(R_3:\Gamma,\Sigma_0\Rightarrow F\).
Apply the smaller-rank lemma on \(F\) to \(R_3\) and \(R_2\),
selecting only the displayed active \(F\) in \(R_2\).  This gives
\(\Gamma,\Sigma_0,\Gamma,\Theta_0\Rightarrow C\).
Collapse \(\Gamma\) once more.  The result is precisely
\(\Gamma,\Pi\Rightarrow C\).
All same-rank calls used the original \(D\) and shorter subtrees of
\(E\); none used the newly constructed \(R_1,R_2,R_3\).

\paragraph{Principal universal quantifier.}
Let \(A=\forall xB\).  Freshen the last eigenparameter of \(D\)
to \(z\), outside all variables of \(E\), its witness term
\(t\), \(\Gamma\), and the chosen matrix representative \(B\).
The left premise is \(D_0:\Gamma\Rightarrow B[z/x]\), and
the right premise is \(E_0:B[t/x],A^{n-1},\Pi\Rightarrow C\).
Cleanup using the original, freshened \(D\) yields
\(R:B[t/x],\Gamma,\Pi\Rightarrow C\).
Apply \([t/z]\) to \(D_0\) using
\ref{thm:p07-lj-alpha-substitution}.  Since \(z\) is not free in
\(\Gamma\), its endsequent is
\(\Gamma\Rightarrow B[t/x]\) as alpha-classes.
The instance \(B[t/x]\) has rank \(\rho(B)<\rho(A)\).
The smaller-rank lemma on this instance and \(R\) gives
\(\Gamma,\Gamma,\Pi\Rightarrow C\); collapse the supplied
contexts.  The substitution's possible increase in term size has no
effect on the rank inequality.

\paragraph{Principal existential quantifier.}
Let \(A=\exists xB\).  The premise of \(D\) is
\(D_0:\Gamma\Rightarrow B[t/x]\).  Freshen the final
eigenparameter of \(E\) to \(z\), outside \(D,t,\Gamma,\Pi,C\)
and the chosen matrix representative.  Its premise is
\(E_0:B[z/x],A^{n-1},\Pi\Rightarrow C\).
Cleanup with the original \(D\) gives
\(R:B[z/x],\Gamma,\Pi\Rightarrow C\).
Substitute \([t/z]\) in this cut-free alpha-proof.  Freshness keeps
\(\Gamma,\Pi,C\) unchanged, and the active formula becomes
\(B[t/x]\) up to alpha-equivalence.  The smaller-rank lemma on
\(B[t/x]\), using \(D_0\), gives
\(\Gamma,\Gamma,\Pi\Rightarrow C\); collapse the contexts.
The proof obtained after cleanup and substitution is used only at
smaller rank, not in a same-rank height argument.

\paragraph{Exhaustion and termination.}
Atoms and bottom have no logical right introductions.  Therefore they
have no remaining principal/principal case: initial sequents and the
left permutations already cover them.  Defined negation is an
implication to bottom, so it is covered by those cases.  The sheet has
no other logical, structural or initial rules once Cut and Mix have
been excluded from the input proofs.

Every same-rank recursive call uses an original immediate premise,
possibly freshened without increasing height, so its height sum is
strictly smaller.  Every other recursive call has smaller formula
rank.  Finite weakening, finite context collapse, and the substitution
constructions terminate independently.  Larger selected multiplicity
in a contraction premise does not alter either measure argument.
Lexicographic induction therefore proves the claimed construction for
all finite input proofs and all selected multiplicities.
\end{proof}
\end{annexitem}

\begin{annexitem}{annex:p7.lj.alpha-identity}{thm:p07-lj-alpha-identity}{Literal alpha-equivalence identities}
\begin{proof}
Induct on the common rank of \(A\) and \(A'\), with the assertion
quantified over ordered pairs of alpha-equivalent literal formulas.
Thus the smaller-rank assertion is available in either direction.
For atoms and bottom, alpha-equivalence is literal equality, so use Id.

For \(A=B\land F\) and \(A'=B'\land F'\), induction gives
\(B\Rightarrow B'\) and \(F\Rightarrow F'\).  Weaken each
to the common antecedent \(B,F\).  Conjunction right gives
\(B,F\Rightarrow B'\land F'\), and conjunction left gives
\(B\land F\Rightarrow B'\land F'\).
For \(A=B\lor F\) and \(A'=B'\lor F'\), inject the inductive
proof \(B\Rightarrow B'\) into the first summand and
\(F\Rightarrow F'\) into the second.  Disjunction left, with
empty common side context, gives the desired conclusion.

For \(A=B\to F\) and \(A'=B'\to F'\), use induction to
obtain \(B'\Rightarrow B\) and \(F\Rightarrow F'\).
The exact split-context implication-left rule yields
\(B\to F,B'\Rightarrow F'\).  Implication right discharges
\(B'\), giving \(B\to F\Rightarrow B'\to F'\).

For \(A=\forall xB\) and \(A'=\forall yB'\), choose a variable
\(z\) fresh for both literal formulas.  The literal instances
\(B[z/x]\) and \(B'[z/y]\) are alpha-equivalent and have smaller
rank.  Induction gives \(B[z/x]\Rightarrow B'[z/y]\).
Universal left with witness \(z\) gives
\(\forall xB\Rightarrow B'[z/y]\), and universal right gives
\(\forall xB\Rightarrow\forall yB'\).  The chosen \(z\)
is not free in the latter conclusion.
For \(A=\exists xB\) and \(A'=\exists yB'\), use the same
inductive instance identity, then existential right on the succedent
and existential left on the antecedent.  The same freshness condition
holds.  Every displayed instance uses the actual total literal
substitution operation; the induction relates its possible literal
alpha-variants.  None of these constructions uses Cut or Mix.
\end{proof}
\end{annexitem}

\begin{annexitem}{annex:p7.lj.literal-realization}{thm:p07-lj-literal-realization}{Realizing the exact literal endsequent}
\begin{proof}
Induct on a cut-free alpha-proof, uniformly over all choices of literal
representatives for the occurrences of its conclusion.  At an identity
leaf the two chosen representatives are alpha-equivalent; apply
\ref{thm:p07-lj-alpha-identity}.  At a bottom-left leaf the representative
of bottom is literally bottom, so the registered initial sequent applies
with all other representatives as chosen.

For WL, choose the premise representatives by deleting the weakened
occurrence from the chosen conclusion.  For CL, choose its two active
premise representatives to be two copies of the exact representative
chosen for the contracted conclusion occurrence.  Realize the premise
by induction and apply the corresponding literal structural rule.

For conjunction left, choose the premise's active formulas to be the
two literal children of the chosen antecedent conjunction.  For
conjunction right, choose the two premise succedents to be the literal
children of the chosen succedent conjunction; copy every common context
representative identically to both premises.  For disjunction left,
choose its two active antecedent formulas as the literal summands of
the chosen conclusion formula and copy the side context identically.
For either disjunction right, choose the premise succedent as the
appropriate literal summand.  The induction hypothesis gives the exact
required premises, to which the respective literal rule applies.

For implication left, retain the recorded split of context occurrences.
If the chosen principal representative is \(B\to F\), choose
\(B\) as the first premise's succedent and \(F\) as the second
premise's active antecedent formula.  Realize both chosen premises
and apply split-context implication left.  For implication right,
choose its premise's discharged assumption and succedent to be the
literal children of the chosen implication, realize that premise,
and discharge precisely that occurrence.

For universal left and existential right, retain the witness term of
the alpha-rule instance.  Form the literal total-substitution instance
of the chosen quantified formula, using its chosen binder name.
Substitution congruence from
\ref{thm:p07-lj-alpha-substitution} puts that literal instance in the
required premise class.  Choose it as the premise representative,
realize the premise by induction, and apply the literal quantifier rule.
For universal right and existential left, use the eigenparameter of
the alpha-rule instance and form the corresponding literal instance
of the chosen quantified formula.  Alpha-equivalent formulas have the
same free variables, so its eigencondition remains true for the chosen
literal conclusion.  The literal instance is in the required premise
class; realize that premise by induction and apply the exact rule.
One may first freshen the parameter in the alpha-premise proof by
\ref{thm:p07-lj-alpha-substitution}; this is unnecessary for validity
but permits a chosen finite avoidance convention.

There are no Cut or Mix nodes to realize.  These cases exhaust the
registered sheet.  In particular, realization at alpha-identity leaves
may increase height.  No height estimate enters this tree induction,
and no realization is used inside the marked-multicut induction.
\end{proof}
\end{annexitem}

\begin{annexitem}{annex:p7.lj.exact-cut-elimination}{thm:p07-lj-exact-cut-elimination}{Elimination of all Cut and Mix nodes}
\begin{proof}
Project the original literal proof to an alpha-proof.  Normalize it by
induction on its original finite tree.  At an initial sequent do nothing.
For an ordinary inference, normalize its premises by the tree induction
hypothesis and reapply the same inference.  Their endsequents are
unchanged, so shared contexts, split contexts and eigenconditions are
unchanged as well.  At a Cut node, its normalized premises contain
neither Cut nor Mix; apply \ref{thm:p07-lj-marked-multicut} with the
single selected cut occurrence.  At a Mix node use its displayed
\(n\geq1\) selected occurrences.  In each case the local lemma
returns a cut-free alpha-proof of the same conclusion.

The original-tree recursion is finite.  Its local multicut operations
terminate by their separate rank/height induction.  No assertion that
one global height measure decreases through all duplications is needed.
Finally apply \ref{thm:p07-lj-literal-realization} to the resulting
cut-free alpha-proof, choosing exactly the literal representatives in
the original endsequent.  The result is a literal derivation in the
exact sheet, with that endsequent and with neither Cut nor Mix.
\end{proof}
\end{annexitem}

\begin{annexitem}{annex:p7.lj.analyticity}{thm:p07-lj-analyticity}{Tracing every formula to the endsequent}
\begin{proof}
For each inference of a cut-free proof, associate each premise formula
occurrence with an occurrence in the conclusion.  A context occurrence
is associated with its corresponding context occurrence.  CL associates
both displayed predecessors with its contracted conclusion occurrence.
WL has no premise occurrence corresponding to its new formula and loses
no existing occurrence.  For \(\land L,\land R,\lor L\) and
\(\lor R_i\), each active constituent is associated with its
principal conjunction or disjunction.  For \(\to L\), associate
the first premise's succedent and the second premise's active antecedent
with the two constituents of the principal implication.  For
\(\to R\), associate the discharged antecedent and the premise
succedent with those two constituents.  Quantifier rules associate
their active matrix instance with the quantified principal formula.

Every formula occurrence therefore has a finite downward path to an
endsequent occurrence.  At each step it is unchanged, a constituent,
or an instance of a quantified matrix.  These operations belong to
the closure in \ref{def:p07-lj-instance-subformula}; iterating them
proves the first-order assertion modulo alpha-equivalence.  In a
propositional derivation no quantifier instance occurs, so the same
argument gives ordinary literal subformulas.  Apply
\ref{thm:p07-lj-exact-cut-elimination} to obtain such a proof from
any derivable sequent.
\end{proof}
\end{annexitem}

\begin{annexitem}{annex:p7.lj.consistency}{thm:p07-lj-consistency}{No derivation of an empty-antecedent bottom sequent}
\begin{proof}
If \(\Rightarrow\bot\) were derivable, cut elimination would give
a cut-free derivation of that same sequent.  It is neither identity
nor bottom left, since both have nonempty antecedents.  No left logical
rule, WL or CL has an empty antecedent in its conclusion: each retains
at least its displayed principal or contracted formula.  No right
logical rule introduces primitive bottom.  There are no right
structural rules, Cut or Mix in the resulting proof.  No last
inference is possible, a contradiction.
\end{proof}
\end{annexitem}

\begin{annexitem}{annex:p7.lj.disjunction-property}{thm:p07-lj-disjunction-property}{Theorem disjunction and the assumption boundary}
\begin{proof}
Eliminate Cut and Mix from a derivation of
\(\Rightarrow B\lor F\).  With empty antecedent, its last step
cannot be an initial sequent, a left logical inference, WL or CL.
The only right rules concluding a disjunction are \(\lor R_1\)
and \(\lor R_2\).  Their respective premises are
\(\Rightarrow B\) and \(\Rightarrow F\), already cut-free.

For the stated limitation with assumptions, fix distinct atoms \(P,Q\).
There is no cut-free proof of \(\Lambda\Rightarrow P\) when every
formula in \(\Lambda\) is either \(Q\) or \(P\lor Q\).
Prove this by height induction.  Identity is impossible because
\(P\) is absent as an antecedent formula, and bottom left is
impossible.  WL or CL would have a smaller premise of the same kind.
The only possible left logical rule is disjunction left; its
\(Q\)-premise again has only \(Q\) and \(P\lor Q\) in the
antecedent, contradicting the induction hypothesis.  No right
introduction concludes the atom \(P\).  These exhaust the cut-free
possibilities.  Taking \(\Lambda\) to contain one \(P\lor Q\)
and using cut elimination rules out \(P\lor Q\Rightarrow P\).
Interchange \(P,Q\) and take the other disjunction premise to rule
out \(P\lor Q\Rightarrow Q\).  Identity still proves
\(P\lor Q\Rightarrow P\lor Q\), establishing the boundary.
\end{proof}
\end{annexitem}

\begin{annexitem}{annex:p7.lj.existence-property}{thm:p07-lj-existence-property}{Witness extraction and grounding}
\begin{proof}
Eliminate Cut and Mix from \(\Rightarrow\exists xB\).  Empty
antecedent excludes all initial and left rules and WL or CL.  The
only right rule with existential succedent is \(\exists R\).
Its premise supplies a term \(t\) and a cut-free proof of
\(\Rightarrow B[t/x]\), as asserted.

Suppose further that \(\exists xB\) is closed and the signature
contains a ground term \(c\).  Let \(\theta\) send every variable
occurring in \(t\) to \(c\), and let \(t_0=t\theta\).
Then \(t_0\) is ground.  Project the extracted proof to alpha-classes
and apply \ref{thm:p07-lj-alpha-substitution} with \(\theta\).
Closedness of \(\exists xB\) means the only possible free variable
of its matrix is its displayed \(x\).  The binding substitution
identities therefore identify the resulting succedent class with
\([B[t_0/x]]_\alpha\).  Literal realization gives a cut-free proof
of exactly \(\Rightarrow B[t_0/x]\).

In the signature of \ref{def:p07-lj-parameter-witness}, the indicated
Id, implication-right and existential-right steps prove the closed
existential formula, but a structural induction on terms shows that
every term is a variable.  No ground term exists.  This proves that
the ground-term hypothesis cannot simply be deleted from the assertion.
\end{proof}
\end{annexitem}

\begin{annexitem}{annex:p7.lj.quantifier-free-conservativity}{thm:p07-lj-quantifier-free-conservativity}{The quantifier-free fragment}
\begin{proof}
Take a cut-free proof of the given quantifier-free endsequent.
Its subformula-instance closure contains no quantifier: taking
constituents and substituting terms cannot create one.  By
\ref{thm:p07-lj-analyticity}, every formula in the proof is therefore
quantifier-free.  Every quantifier rule has a quantified principal
formula in its conclusion, so no quantifier rule occurs.  In the
remaining inference rules, each active premise formula is a literal
constituent of a principal conclusion formula, and all other premise
formulas are unchanged context occurrences.  Following these
associations to the endsequent gives literal subformulas.  The
remaining rule instances operate only on \(\bot,\land,\lor,\to\)
and leave each first-order atom intact.  Reading each such atom as a
propositional atom thus gives exactly the asserted fragment derivation.
\end{proof}
\end{annexitem}
